\chapter{Moduli Spaces In Supersymmetric Quantum Field Theory}
\label{ch:susy-moduli-spaces}

\ConventionsMixed

This chapter defines supersymmetric moduli spaces as spaces of supersymmetric
vacua together with their low-energy QFT data.  The set of classical solutions
is only the first coordinate chart; the physical object includes quotienting,
singular strata, massless fields, chiral rings, metrics, and quantum
corrections.

\section{Vacuum Equations}

For four-dimensional \(\mathcal N=1\) chiral multiplets \(\Phi^i\) with
superpotential \(W(\Phi)\) and gauge group \(G\), a classical supersymmetric
vacuum obeys
\[
  F_i=\frac{\partial W}{\partial\Phi^i}=0,
  \qquad
  \mu^a=0,
\]
where \(\mu:\mathcal X\to\mathfrak g^\ast\) is the moment map for the
K\"ahler target \(\mathcal X\) of scalar fields.  The classical vacuum space is
the K\"ahler quotient
\[
  \mathcal M_{\mathrm{cl}}
  =
  \{F=0,\ \mu=0\}/G .
\]
Equivalently, under the required stability hypotheses for the complexified gauge
action, it is the holomorphic quotient
\[
  \{F=0\}/G_\mathbb C .
\]
The second description identifies holomorphic functions on
\(\mathcal M_{\mathrm{cl}}\) with gauge-invariant chiral polynomials modulo
the ideal generated by the \(F\)-term equations.

The last sentence hides several hypotheses, so we spell out the quotient
datum used in this volume.  Let \(V\) be the finite-dimensional Hermitian
representation of \(G\) carried by the scalar fields, let \(G_\mathbb C\) be a
reductive complexification, and let \(W\in\mathbb C[V]^{G_\mathbb C}\) be a
polynomial superpotential.  Write
\[
  Z_F=\{v\in V\mid dW(v)=0\},\qquad
  I_F=(\partial_1W,\ldots,\partial_nW)\subset\mathbb C[V].
\]
The affine chiral quotient is
\[
  \mathcal M_{\mathrm{ch},\mathrm{cl}}
  =
  \operatorname{Spec}
  \left((\mathbb C[V]/I_F)^{G_\mathbb C}\right)_{\mathrm{red}},
  \label{eq:affine-chiral-quotient}
\]
where the subscript records that the displayed affine variety forgets
nilpotents unless they are kept as derived or scheme-theoretic data.  The
symplectic vacuum quotient is
\[
  \mathcal M_{\mu,\mathrm{cl}}
  =
  (Z_F\cap\mu^{-1}(0))/G .
  \label{eq:symplectic-vacuum-quotient}
\]

\begin{proposition}[Classical quotient and chiral functions]
\label{prop:susy-classical-quotient-chiral-ring}
Assume the Kempf--Ness comparison for the triple
\((V,G,G_\mathbb C)\) on the \(G_\mathbb C\)-invariant closed subset
\(Z_F\): every polystable \(G_\mathbb C\)-orbit in \(Z_F\) meets
\(\mu^{-1}(0)\), and two points of \(Z_F\cap\mu^{-1}(0)\) lie in the same
closed \(G_\mathbb C\)-orbit exactly when they lie in the same \(G\)-orbit.
Then \(\mathcal M_{\mu,\mathrm{cl}}\) is naturally identified with the
polystable affine quotient \(\mathcal M_{\mathrm{ch},\mathrm{cl}}\), and its
holomorphic coordinate ring is
\[
  \mathbb C[\mathcal M_{\mathrm{ch},\mathrm{cl}}]
  =
  \left((\mathbb C[V]/I_F)^{G_\mathbb C}\right)_{\mathrm{red}} .
\]
Thus the notation \(\{F=0\}/G_\mathbb C\) in this monograph means the affine
quotient in \eqref{eq:affine-chiral-quotient}, together with the stated
stability hypothesis, not the naive set of all complexified orbits.
\end{proposition}

\begin{proof}
Because \(W\) is \(G_\mathbb C\)-invariant, the differential \(dW\) is
equivariant and the ideal \(I_F\) is \(G_\mathbb C\)-stable.  Hence every
class in \((\mathbb C[V]/I_F)^{G_\mathbb C}\) is represented by a
gauge-invariant holomorphic polynomial on \(Z_F\), and restriction to
\(Z_F\cap\mu^{-1}(0)\) gives a holomorphic function constant on \(G\)-orbits.
The assumed Kempf--Ness statement identifies \(G\)-orbits in
\(Z_F\cap\mu^{-1}(0)\) with closed \(G_\mathbb C\)-orbits in the affine
quotient.  Invariant polynomials separate closed orbits in an affine
reductive quotient, so the descended functions are precisely the regular
holomorphic functions on the reduced quotient.  This proves the asserted
coordinate-ring identification.  The final warning follows because nonclosed
complexified orbits can have the same closure and therefore the same values of
all invariant polynomials; the affine quotient collapses them to the
polystable representative.
\end{proof}

\begin{example}[Rank-one abelian quotient]
\label{ex:susy-rank-one-abelian-quotient}
Let \(G=U(1)\) act on \(V=\mathbb C^2\) by
\[
  (x,y)\longmapsto(\ee^{\ii\alpha}x,\ee^{-\ii\alpha}y),
  \qquad W=0,
\]
with moment map \(\mu=|x|^2-|y|^2\).  The invariant ring of the
complexified action \(t\cdot(x,y)=(tx,t^{-1}y)\) is
\[
  \mathbb C[x,y]^{\mathbb C^\ast}=\mathbb C[u],
  \qquad u=xy .
\]
Indeed the monomial \(x^ay^b\) has \(U(1)\) charge \(a-b\), so invariance
forces \(a=b\) and the monomial is \(u^a\).  On the symplectic side
\(\mu=0\) gives \(|x|=|y|\).  The map \((x,y)\mapsto xy\) sends
\(\mu^{-1}(0)/U(1)\) bijectively to \(\mathbb C\): for
\(u=\rho\ee^{\ii\theta}\), with \(\rho\ge0\), choose
\[
  x=\sqrt{\rho},\qquad y=\sqrt{\rho}\,\ee^{\ii\theta},
\]
and changing the phase split between \(x\) and \(y\) is precisely the
\(U(1)\) quotient.  Thus this basic vacuum quotient is the complex line with
coordinate \(u\), not a two-complex-dimensional space of unquotiented scalar
expectation values.
\end{example}

\section{Quantum Moduli Space Datum}

A quantum moduli space is a tuple
\[
  (\mathcal M,\mathcal H_p,\mathcal A_p,\mathcal G_p)_{p\in\mathcal M},
\]
where \(\mathcal M\) is the parameter space of supersymmetric vacua,
\(\mathcal H_p\) is the Hilbert space of the theory in vacuum \(p\),
\(\mathcal A_p\) is the algebra of local operators in the low-energy theory,
and \(\mathcal G_p\) records background responses and anomaly data.  On a
smooth branch with massless scalar coordinates \(z^i\), the two-derivative
effective action determines a metric
\[
  ds^2=g_{i\bar j}(z,\bar z)\,dz^i d\bar z^{\bar j}
\]
on the branch.  The metric is part of the low-energy QFT data; the chiral
coordinate ring determines only holomorphic functions.

\section{Chiral Rings And Holomorphic Functions}

Let \(\mathcal R_{\mathrm{ch}}\) be the chiral ring of gauge-invariant local
operators modulo \(\bar Q\)-exact operators.  In a vacuum \(p\), expectation
values define a homomorphism
\[
  \operatorname{ev}_p:\mathcal R_{\mathrm{ch}}\longrightarrow\mathbb C .
\]
If chiral operators separate vacua on a branch and have no nilpotent
relations in the reduced branch coordinate ring, the collection of
\(\operatorname{ev}_p\) identifies the reduced holomorphic coordinate ring
with a quotient of \(\mathcal R_{\mathrm{ch}}\).  Nilpotents, accidental
massless fields, and singular strata must be retained when they affect
operator products or low-energy dynamics.

\section{N One SQCD Branches}

For \(SU(N_c)\) SQCD with \(N_f\) flavors, the gauge-invariant chiral
coordinates are mesons and baryons as defined in
Chapter~\ref{ch:susy-four-dimensional-n1-gauge-dynamics}.  The branch data
depend on \(N_f\):
\[
\begin{array}{c|c}
N_f<N_c & \text{ADS superpotential produces a runaway direction}\\
N_f=N_c & \det M-B\widetilde B=\Lambda_h^{2N_c}\\
N_f=N_c+1 & \text{confining chiral coordinates with superpotential constraint}\\
\end{array}
\]
Each entry is a Wilsonian chiral-coordinate statement.  The metric on the
branch and the massive spectrum require additional low-energy data.

\section{N Two Coulomb, Higgs, And Mixed Branches}

For \(\mathcal N=2\) theories, Coulomb-branch coordinates come from vector
multiplet scalars and have special K\"ahler geometry.  Higgs branches come from
hypermultiplet scalars and are hyperk\"ahler quotients:
\[
  \mathcal M_H=\mu_\mathbb R^{-1}(0)\cap\mu_\mathbb C^{-1}(0)/G .
\]
Mixed branches have both vector and hypermultiplet massless fields after
solving compatibility conditions between vector scalar expectation values
and hypermultiplet expectation values.  The Coulomb-branch metric may receive
quantum corrections encoded by Seiberg-Witten periods.  In many Lagrangian
examples the Higgs-branch hyperk\"ahler metric is protected, but this protection
is a statement about the supersymmetric Wilsonian scheme and the branch
coordinates, not a statement about arbitrary coordinate choices.

\section{Singularities}

A point \(p\in\mathcal M\) is singular when the low-energy description changes
there: additional states can become massless, the effective metric can
degenerate, or the coordinate quotient can have an orbifold or more severe
singularity.  In the \(\mathcal N=2\) pure gauge-algebra \(\mathfrak{su}(2)\)
example, the singularities \(u=\pm\Lambda^2\) are interpreted by vanishing
cycles of the Seiberg-Witten curve and the appearance of light monopole or
dyon hypermultiplets.  The correct local model is the Abelian theory with
those light fields included.

\section{Effective Theory On A Branch}

Let \(p\) lie on a smooth branch and let \(z^i\) be local coordinates.  The
low-energy effective action has the derivative expansion
\[
  S_{\mathrm{eff}}
  =
  \int d^dx\,
  g_{i\bar j}(z,\bar z)\partial_\mu z^i\partial^\mu\bar z^{\bar j}
  +\text{fermions}
  +\text{higher-derivative terms}.
\]
The branch approximation is valid for momenta small compared with the masses
of fields integrated out and away from loci where additional fields become
light.  A complete moduli-space statement therefore includes its domain of
validity inside the full QFT.

\begin{openproblem}[Moduli-space construction]
\label{op:susy-nonperturbative-moduli-space-construction}
Construct supersymmetric moduli spaces directly from regulated continuum QFT:
the vacuum sectors, chiral-coordinate ring, low-energy Hilbert spaces,
branch metrics, singular local models, and background anomaly responses.
\end{openproblem}
